\documentclass{article}

\usepackage{tabularx}
\usepackage{booktabs}

\title{Reflection and Traceability Report on \progname}

\author{\authname}

\date{}

\input{../Comments}
\input{../Common}

\begin{document}

\maketitle

\plt{Reflection is an important component of getting the full benefits from a
learning experience.  Besides the intrinsic benefits of reflection, this
document will be used to help the TAs grade how well your team responded to
feedback.  Therefore, traceability between Revision 0 and Revision 1 is and
important part of the reflection exercise.  In addition, several CEAB (Canadian
Engineering Accreditation Board) Learning Outcomes (LOs) will be assessed based
on your reflections.}

\section{Changes in Response to Feedback}

\plt{Summarize the changes made over the course of the project in response to
feedback from TAs, the instructor, teammates, other teams, the project
supervisor (if present), and from user testers.}

\plt{For those teams with an external supervisor, please highlight how the feedback 
from the supervisor shaped your project.  In particular, you should highlight the 
supervisor's response to your Rev 0 demonstration to them.}

\plt{Version control can make the summary relatively easy, if you used issues
and meaningful commits.  If you feedback is in an issue, and you responded in
the issue tracker, you can point to the issue as part of explaining your
changes.  If addressing the issue required changes to code or documentation, you
can point to the specific commit that made the changes.  Although the links are
helpful for the details, you should include a label for each item of feedback so
that the reader has an idea of what each item is about without the need to click
on everything to find out.}

\plt{If you were not organized with your commits, traceability between feedback
and commits will not be feasible to capture after the fact.  You will instead
need to spend time writing down a summary of the changes made in response to
each item of feedback.}

\plt{You should address EVERY item of feedback.  A table or itemized list is
recommended.  You should record every item of feedback, along with the source of
that feedback and the change you made in response to that feedback.  The
response can be a change to your documentation, code, or development process.
The response can also be the reason why no changes were made in response to the
feedback.  To make this information manageable, you will record the feedback and
response separately for each deliverable in the sections that follow.}

\plt{If the feedback is general or incomplete, the TA (or instructor) will not
be able to grade your response to feedback.  In that case your grade on this
document, and likely the Revision 1 versions of the other documents will be
low.} 
\subsection{SRS and Hazard Analysis}

Over the course of the project, we made targeted updates to the SRS and Hazard
Analysis in response to TA feedback, advisor guidance, and internal review. The
main themes were: (i) improving requirement precision and testability, (ii)
clarifying system boundaries and privacy assumptions for a local-first assistive
tool, and (iii) strengthening safety language around command execution and user
confirmation.

\begin{itemize}

\item \textbf{Clarified system boundary and privacy assumptions}

\textbf{Source:} TA/advisor feedback during Revision 0 review and demo discussions.

\textbf{Feedback:} The boundary between what runs locally vs.\ what relies on
external components was not explicit enough, which made feasibility and privacy
risk harder to assess.

\textbf{Response:} We revised the SRS to state the local-first intent more
clearly, identify external dependencies where applicable, and tighten definitions
around what user data exists, where it can be stored, and what is not retained by
default.

\item \textbf{Improved requirement specificity and testability}

\textbf{Source:} TA feedback on requirement quality.

\textbf{Feedback:} Some requirements were too high-level or combined multiple
behaviours, reducing testability and traceability.

\textbf{Response:} We reworded several requirements to be more measurable (clear
triggers, observable outputs, and success criteria) and split bundled
requirements so each one maps cleanly to design elements and V\&V tests.

\item \textbf{Strengthened traceability between SRS/HA and the design docs}

\textbf{Source:} TA feedback emphasizing traceability, including the design
documentation review
(\href{https://github.com/speech-buddies/VoiceBridge/issues/235}{TA Review \#235}).

\textbf{Feedback:} The design documentation needed clearer, more consistent
module responsibilities and mappings, and the project artifacts needed to remain
traceable across documents.

\textbf{Response:} While \#235 primarily drove updates to the design documents,
addressing it required us to cross-check the design against the SRS and Hazard
Analysis. We verified and tightened SRS-to-module and hazard/mitigation-to-module
traceability to reduce drift and ensure that safety and privacy assumptions in
the SRS/HA were reflected in the architecture.

\item \textbf{Expanded hazards related to unintended command execution}

\textbf{Source:} TA feedback and internal risk review after implementing command
execution flows.

\textbf{Feedback:} The hazard analysis under-emphasized risks from
misrecognition, ambiguous intent, and executing harmful browser actions.

\textbf{Response:} We strengthened hazards and mitigations around incorrect
transcriptions and ambiguous commands leading to unintended actions. Mitigations
were updated to emphasize explicit confirmation for higher-risk actions, clear
user feedback during execution, and safe fallback behaviour when commands are
out-of-scope or low-confidence.

\item \textbf{Made privacy and data-retention hazards more concrete}

\textbf{Source:} TA/advisor feedback on privacy risk articulation.

\textbf{Feedback:} Privacy hazards were present but too abstract, making it
unclear what data might be logged or retained and what controls exist.

\textbf{Response:} We clarified which artifacts (e.g., transcripts, recordings,
logs, and user settings) might exist and refined mitigations to emphasize
minimizing stored sensitive data, limiting log exposure, and treating stored
artifacts as sensitive by default.

\item \textbf{Addressed accessibility-related failure modes as safety/usability risks}

\textbf{Source:} TA feedback and usability-oriented review.

\textbf{Feedback:} Since the target users include individuals with dysarthric
speech, the analysis should consider failure modes such as unclear prompts,
missing confirmations, or confusing system state.

\textbf{Response:} We added/refined hazards capturing user confusion and loss of
control due to insufficient feedback, and we connected mitigations to clearer
state indications and consistent confirm/cancel interactions.

\end{itemize}

\subsection{Design and Design Documentation}

Most of our feedback-driven changes were in the design documentation, where we
focused on improving clarity, consistency, and alignment with the implemented
system.

\begin{itemize}

\item \textbf{Improved module boundaries and responsibility descriptions}

\textbf{Source:} TA feedback on ambiguity/overlap between components.

\textbf{Feedback:} Some responsibilities were described in overlapping terms
(e.g., session flow vs.\ command interpretation vs.\ browser execution), making
it hard to understand component ownership.

\textbf{Response:} We revised the Module Guide's module descriptions to better
separate responsibilities (UI/feedback vs.\ transcription vs.\ command synthesis
vs.\ execution/orchestration vs.\ storage/security). This made the architecture
easier to understand and reduced ambiguity for implementation and testing.

\item \textbf{Updated documentation to match the final implementation}

\textbf{Source:} TA/advisor feedback plus internal consistency checks.

\textbf{Feedback:} Some design text reflected earlier architecture drafts rather
than the final code structure, creating ``documentation drift.''

\textbf{Response:} We updated module descriptions and terminology so they
accurately reflect the final architecture and how the system is implemented
(including details of the audio capture flow and the model tuning workflow).
This improved the MG's usefulness as a maintenance reference.

\item \textbf{Strengthened traceability matrices and requirement coverage}

\textbf{Source:} TA feedback emphasizing traceability and completeness.

\textbf{Feedback:} The design needed clearer evidence that requirements
(functional, usability, performance, security, etc.) were mapped to concrete
parts of the system.

\textbf{Response:} We revised and validated the traceability matrices so each
requirement category maps to the relevant modules, and we checked for
gaps/duplicates. This also helped ensure the V\&V work could be planned directly
against the design.

\item \textbf{Clarified inter-module communication and error handling}

\textbf{Source:} Advisor/TA feedback requesting clearer interaction details.

\textbf{Feedback:} The documents needed clearer description of how modules
exchange information and how errors propagate to user-facing feedback.

\textbf{Response:} We expanded the design documentation to describe the
message/data flow at a high level (what data types are passed, where session
state is tracked, and how failures are turned into user-facing feedback and
logs). This made the design more actionable and easier to verify.

\item \textbf{Refined UI and accessibility documentation based on feedback}

\textbf{Source:} TA feedback and informal user-focused review.

\textbf{Feedback:} The UI story needed clearer user-state feedback (listening,
processing, confirm/cancel) and accessibility justification.

\textbf{Response:} We improved UI-related sections and figures/descriptions to
emphasize predictable state transitions, clear confirmations, and accessible
feedback. These changes were motivated by the needs of users with dysarthric
speech, where system confidence and clarity directly affect usability.

\end{itemize}

\subsection{VnV Plan and Report}

\section{Extras}

\plt{Summarize the extras that were tackled by this project.  Extras
can include usability testing, code walkthroughs, user documentation, formal
proof, GenderMag personas, Design Thinking, etc.  Extras should have already
been approved by the course instructor as included in your problem statement.}

\section{Design Iteration (LO11 (PrototypeIterate))}

\plt{Explain how you arrived at your final design and implementation.  How did
the design evolve from the first version to the final version?} 

\plt{Don't just say what you changed, say why you changed it.  The needs of the
client should be part of the explanation.  For example, if you made changes in
response to usability testing, explain what the testing found and what changes
it led to.}

\subsection{Shortcut Manager Iteration}

One important design iteration in VoiceBridge was the development of the
\texttt{ShortcutManager}. In the earlier versions of the system, commands were
handled one at a time with no persistent way to record, save, and replay common
action sequences. While this was enough to demonstrate basic browser control, it
did not match the client need for efficiency and accessibility in repeated
everyday tasks. For a user with impaired speech, repeatedly issuing the same
sequence of voice commands can increase fatigue and reduce the overall usability
of the system. Because of this, we introduced shortcuts as a way to reduce the
number of spoken interactions required for frequent tasks.

Our first idea was relatively simple: store a shortcut as a sequence of recorded
commands and replay them later. However, once we integrated this into the live
system, we found that shortcut behavior could conflict with the runtime state
machine. For example, the user could try to begin a shortcut recording while the
system was still processing another command, or multiple parts of the backend
could attempt to modify shortcut state at the same time. This was especially
important because VoiceBridge is a real-time system with asynchronous events,
including speech capture, transcription, intent resolution, and browser
execution. As a result, we revised the design so that shortcut recording became
an explicitly managed subsystem rather than just an added list of commands.

The final design introduced a dedicated \texttt{ShortcutManager} module with a
clear internal state, persistent storage, and locking for safe concurrent
access. Instead of scattering shortcut logic across the codebase, the manager
became responsible for starting and stopping recordings, storing recorded
commands, validating shortcut names, and loading/saving shortcuts from disk.
This design was chosen to improve maintainability and prevent race conditions in
cases where multiple backend events happen close together. The use of a lock was
an important iteration because testing showed that shortcut actions could overlap
with other command-processing flows in unpredictable ways if the state was not
protected.

We also changed how shortcuts were represented. In earlier thinking, shortcuts
were treated mainly as temporary in-memory recordings for the active session.
This was limiting, because a user would lose their shortcuts between runs of the
application. Since one of the client goals was to make the system practical for
daily use, we moved to persistent storage using a local JSON-based approach.
This allowed shortcuts to survive application restarts and made the feature more
realistic as an assistive tool rather than just a demo feature. The trade-off
was that we now needed validation and structured file handling, but this was
worth it because it supported a more reliable user experience.

Another design change was motivated by usability. Initially, shortcut support was
focused only on storing the raw sequence of actions. As the design matured, we
recognized that the feature also needed clear start/stop control and user
feedback so that recording would feel predictable. Without this, a user could be
unsure whether the system was actively recording commands into a shortcut or just
executing them normally. This was especially important given our client focus on
accessibility and confidence in system behavior. We therefore refined the design
so that shortcut recording had explicit transitions, clear prompts, and tighter
integration with the system state flow.

Overall, the \texttt{ShortcutManager} evolved from a simple convenience idea
into a more structured subsystem designed around accessibility, predictability,
and safe execution. The main reason for these changes was not just technical
cleanup, but the need to better support users who benefit from reducing repeated
speech effort while still being able to trust what the system is doing.

\section{Design Decisions (LO12)}

\plt{Reflect and justify your design decisions.  How did limitations,
 assumptions, and constraints influence your decisions?  Discuss each of these
 separately.}

\section{Economic Considerations (LO23)}

\plt{Is there a market for your product? What would be involved in marketing your 
product? What is your estimate of the cost to produce a version that you could 
sell?  What would you charge for your product?  How many units would you have to 
sell to make money? If your product isn't something that would be sold, like an 
open source project, how would you go about attracting users?  How many potential 
users currently exist?}

\section{Reflection on Project Management (LO24)}

\subsection{How Does Your Project Management Compare to Your Development Plan}

Our project management aligned well with our Development Plan. We maintained a
consistent meeting rhythm to coordinate priorities and remove blockers, and we
used asynchronous communication between meetings to keep work moving. We used
GitHub Issues and pull requests to plan tasks, record feedback, and track progress
across deliverables. Team roles were clear and helped us work efficiently, while
still allowing collaboration across areas when needed. We also used the
technologies we planned on using (GitHub for version control and coordination,
and LaTeX for documentation), and we kept our workflow structured around
reviewable changes and traceability.

\subsection{What Went Well?}

\begin{itemize}
  \item \textbf{Issue tracking and traceability:} Capturing peer and TA feedback
  as GitHub Issues made it easy to plan changes and clearly demonstrate
  traceability from comments to updates.
  \item \textbf{Clear ownership with collaboration:} Assigning primary owners for
  major deliverables improved consistency and helped each document progress
  steadily, while collaboration ensured quality and shared understanding.
  \item \textbf{Review and iteration:} Regular internal review cycles improved
  clarity, consistency, and alignment across the SRS/HA, design documentation,
  and V\&V artifacts.
  \item \textbf{Coordination across implementation and documentation:} Keeping
  documentation connected to the implementation helped ensure design decisions
  were well-justified and supported by the final system.
\end{itemize}

\subsection{What Went Wrong?}

\begin{itemize}
  \item \textbf{Continuous alignment effort:} Maintaining alignment between
  requirements, design, and V\&V artifacts required ongoing attention, and the
  experience reinforced the value of regularly revisiting traceability links as
  the system evolves.
  \item \textbf{Nonfunctional evaluation planning:} Defining measurable
  nonfunctional tests and presenting quantitative results clearly was a valuable
  learning area that strengthened the quality of our V\&V work.
  \item \textbf{Estimating integration work:} Coordinating integration across
  components highlighted the importance of explicitly planning interface details
  and test hooks early, which improved the final system robustness.
\end{itemize}

\subsection{What Would you Do Differently Next Time?}

\begin{itemize}
  \item \textbf{Make traceability updates routine:} Update traceability matrices
  on a regular schedule so cross-document consistency is maintained continuously.
  \item \textbf{Define acceptance criteria for issues:} Add short, explicit
  acceptance criteria to issues to keep scope clear and improve verification.
  \item \textbf{Plan earlier nonfunctional baselines:} Establish performance and
  usability baselines earlier so evaluation is repeatable and comparisons across
  iterations are straightforward.
\end{itemize}

\section{Ethics and Equity (LO18)}

Equity and universal design were central to VoiceBridge because the system is
intended to support individuals with dysarthric speech. We designed for clarity,
predictability, and user control by emphasizing clear system-state feedback
(listening, processing, executing), confirmation for higher-risk actions, and
safe handling of low-confidence or out-of-scope commands. We also treated privacy
as part of equitable design by documenting data assumptions and encouraging
local-first handling where feasible, so users can understand and trust how their
information is processed. Across the project, we aimed to reduce unnecessary
cognitive load and promote confidence by keeping interactions consistent and
transparent.

\section{Reflection on Capstone}

\subsection{Which Courses Were Relevant}

\begin{itemize}
  \item \textbf{SFWRENG 4HC3 (Human-Computer Interaction):} Informed our focus on
  usable, predictable interactions, user feedback, and accessibility-oriented
  design decisions.
  \item \textbf{SFWRENG 3A04 (Software Architecture):} Supported our approach to
  module decomposition, clear responsibility assignment, and documenting
  interfaces and system structure.
  \item \textbf{SFWRENG 4AL3 (Applied Machine Learning):} Helped us reason about
  ASR-related considerations such as evaluation, error modes, and performance
  tradeoffs in practical settings.
\end{itemize}

\subsection{Knowledge/Skills Outside of Courses}

\begin{itemize}
  \item \textbf{End-to-end ASR pipeline integration:} Building and integrating a
  real-time flow from audio capture to transcription to command execution.
  \item \textbf{Practical evaluation:} Designing repeatable procedures and
  metrics (e.g., latency and error characterization) that reflect real usage.
  \item \textbf{Privacy-aware engineering:} Translating privacy expectations into
  concrete requirements, hazards, and mitigations for an assistive tool.
  \item \textbf{Documentation and maintenance discipline:} Keeping requirements,
  design, and V\&V artifacts consistent through traceability as the project
  evolved.
  \item \textbf{Collaborative workflow:} Applying issue tracking, pull requests,
  and structured reviews to coordinate work and justify changes.
\end{itemize}

\end{document}